\subsection{A faithful geometric realization of the full specialized source algebra}

\begin{proposition}[P-F73: all twenty-one source dimensions survive and act geometrically]
\label{prop:nonstandardfullcorrespondence}
Fix $q=3^e$, $e\ge1$.  Let $\mathcal B_q$ be the rational
specialization of the actual convenient word form
in (P-F52ai), with its original generators $Q_1,Q_2$.
Put
\[
 \begin{split}
 f&=q^5+q+q^{-1}+q^{-5},\\
 g_3&=q^4+2+q^{-4},\quad
 g_4=q^6+2+q^{-6},\quad
 g_5=q^{10}+2+q^{-10},\\
 K&=\mathbb Q(\sqrt{A(q)}),\qquad
 \gamma=S-\alpha(q)/F\in K,\\
 M_1(h)&=\begin{pmatrix}f&h\\0&0\end{pmatrix},\qquad
 M_2(h)=\begin{pmatrix}0&0\\1&f\end{pmatrix}.
 \end{split}                                                   \tag{P-F73a}
\]
Here $\alpha(q)$, $S$, and $F$ have exactly the source
normalizations fixed in P-F70.  There is an explicit
algebra isomorphism
\[
 \begin{split}
 \Phi:\mathcal B_q&\xrightarrow{\ \sim\ }
       \mathbb Q\oplus M_2(\mathbb Q)^3\oplus M_2(K),\\
 Q_i&\longmapsto
       \bigl(0,M_i(g_3),M_i(g_4),M_i(g_5),M_i(\gamma)\bigr),
                    \qquad i=1,2.
 \end{split}                                                   \tag{P-F73b}
\]
The specialization of the source's actual tensor action
on $X_q^{\otimes3}$ is faithful on $\mathcal B_q$.
Thus (P-F73b) describes that full specialized image,
not only an abstract block selected from it.

Let $\mathcal A/\mathbb F_Q$, $Q=q^{116}$, be the
surface from P-F69, with the $K$-embedding fixed in P-F70.
Let $\mathcal E/\mathbb F_Q$ be the elliptic curve
\[
                         y^2=x^3-x.
                                                               \tag{P-F73c}
\]
For the eleven-dimensional abelian variety
\[
                    \mathcal X=\mathcal A^2\times\mathcal E^7
                                                               \tag{P-F73d}
\]
the isomorphism (P-F73b) gives a faithful unital algebra map
\[
             \mathcal B_q\hookrightarrow
                         \operatorname{End}_{\mathbb F_Q}^0(\mathcal X).
                                                               \tag{P-F73e}
\]
Using blockwise transposition before contravariant pullback
gives a faithful left action of the same source algebra
on the pure weight-one space
$H^1(\mathcal X_{\overline{\mathbb F}_Q},\mathbb Q_\ell)$,
of dimension $22$, for every $\ell\ne3$.
It preserves the specified generators, units, ordered
words, all relations, and the actual scalar rescaling
$Q_i\mapsto FQ_i$.
\end{proposition}

\begin{proof}
\par\medskip\noindent\textbf{1. The original presentation and its specialization.}
The exact source is GCT VII, Section~7.1.1, printed
pp.~18--19, the length-eleven relation in Figure~2,
printed p.~31, and its image under $1\leftrightarrow2$:
\url{https://arxiv.org/pdf/0709.0749}.
The source explicitly gives the $21$ alternating words
of lengths zero through ten as a basis.  Section~7.1.2,
pp.~19--20, also records the generic one-dimensional
and five two-dimensional block sizes.  The proof here
uses the presentation and word basis, not the
quarantined triangular matrix display on p.~20.

Over the formal Laurent ring $R_0=\mathbb Z[v,v^{-1}]$
the relations are $Q_i^2=fQ_i$ and
\[
                    r(G)Q_1=r(G)Q_2=0,
 \qquad G=f^2-(Q_1-Q_2)^2,
                                                               \tag{P-F73f}
\]
where P-F52b--f3 proves the centrality of $G$ and the
complete coefficient identity
\[
 \begin{split}
 r(T)&=(T^2-L_0T+N_0)(T-g_3)(T-g_4)(T-g_5),\\
 L_0&=3(v^8+v^{-8})+4(v^4+v^{-4})
                               -2(v^2+v^{-2})+10,\\
 N_0&=\frac{(v^4-v^2+1)^2
          (v^{12}+v^{10}+2v^8+2v^4+v^2+1)^2}{v^{16}}.
 \end{split}                                                   \tag{P-F73g}
\]
All scalars in this paragraph are evaluated at the
formal parameter $v$ until specialization.
The monic quadratic and length-eleven reductions
strictly shorten words and have coefficients in
$R_0$.  Hence the $21$ words span the word form over
$R_0$; their cited independence over $\mathbb Q(v)$
makes them an $R_0$-basis.  This is also the integral
word-form argument already established at (P-F52ai).
Specializing $v$ to the nonzero rational number $q$
therefore leaves a $21$-dimensional rational algebra
with precisely those relations and basis vectors.

\par\medskip\noindent\textbf{2. Noncollision and the complete algebra decomposition.}
At $q=3^e$, the discriminant of the quadratic in
(P-F73g) is
\[
             L_0^2-4N_0=\frac{D(q)}{q^{40}}
                       =q^{-16}(1+q^8)^2A(q).
                                                               \tag{P-F73h}
\]
Because $A(q)\equiv2\pmod3$, this is nonsquare in
$\mathbb Q$, and that quadratic is irreducible.
Its two roots are $g_1,g_2$ in (P-F52f1), and
$\gamma=S-\alpha/F$ is exactly one of them, with the
specified $K$-identification from (P-F70h).
In particular neither root is rational, so neither
is $0$, $f^2$, or any $g_j$, $3\le j\le5$.

For $q>1$, the function $x+x^{-1}$ is strictly
increasing for $x>1$, and hence
\[
                        0<g_3<g_4<g_5<f^2.
                                                               \tag{P-F73i}
\]
For the last inequality the exact expansion is
\[
 f^2=q^{10}+q^{-10}+q^2+q^{-2}
             +2(q^6+q^{-6})+2(q^4+q^{-4})+4.
                                                               \tag{P-F73j}
\]
Subtracting $g_5$ leaves a Laurent polynomial with
strictly positive coefficients.  These observations
prove that the five rational irreducible factors of
\[
                      \mathcal P(T)=(T-f^2)r(T)
                                                               \tag{P-F73k}
\]
are pairwise coprime at the specified specialization.

For any scalar $h$ in a commutative field, direct
multiplication gives
\[
 M_i(h)^2=fM_i(h),\qquad
             f^2I-(M_1(h)-M_2(h))^2=hI.
                                                               \tag{P-F73l}
\]
Thus the four matrix components in (P-F73b) satisfy
both source relations (P-F73f), because their $h$
values are roots of $r$.  The scalar component has
$Q_1=Q_2=0$ and also satisfies both relations.  This
proves that (P-F73b) defines a homomorphism from the
actual presentation.  The scalar component agrees
with the actual source's $T_0$, for which p.~20
explicitly states $Q_j u_0=0$ for $j=1,2$.

Each matrix component is the full indicated matrix
algebra.  Indeed its image contains $hI$, and
\[
 \begin{split}
 fM_1-M_1M_2&=(f^2-h)E_{11},\\
 fM_1M_2-hM_1&=h(f^2-h)E_{12},\\
 E_{22}&=I-E_{11},\qquad E_{21}=M_2-fE_{22}.
 \end{split}                                                   \tag{P-F73m}
\]
The denominators $h(f^2-h)$ are nonzero by the
preceding paragraph.  In the quadratic component
$\mathbb Q[h]=K$, so their inverses belong to that
field and are polynomial expressions in $h$.

Finally $\Phi(G)$ has scalar values $f^2,g_3,g_4,g_5$
on the rational components and $\gamma$ on the
quadratic one.  The Chinese remainder theorem for
the five coprime factors of (P-F73k) gives polynomials
in $G$ whose images are the five component identities.
They isolate the scalar component and each of the
full matrix images just proved.  Therefore $\Phi$
is surjective onto their direct product.  The latter
has rational dimension
$1+3\cdot4+4[K:\mathbb Q]=21$, equal to the dimension
of its source.  Thus $\Phi$ is an isomorphism.

\par\medskip\noindent\textbf{3. No loss in the actual tensor specialization.}
We verify this separately; abstract semisimplicity
alone would not establish faithfulness of a particular
specialized representation.
The source's actual two-tensor rescaled projection
matrix $\mathsf P$ in Figure~1, printed p.~30, has
Laurent-polynomial entries in the formal parameter:
every displayed denominator is a power of $v$.
Its complementary matrix is $\mathsf Q=fI-\mathsf P$.
Consequently the actual matrices
$Q_1=\mathsf Q\otimes I$, $Q_2=I\otimes\mathsf Q$,
and $G$ on the $64$-dimensional tensor cube have
entries in $\mathbb Q[v,v^{-1}]$ and specialize
regularly at every nonzero $v$.
These are the matrices whose actual positive tensor
form was checked in P-F52k--l; no change to the
printed two-dimensional block matrices is used.

The same decomposition argument in step~2 applies
over $\mathbb Q(v)$.  The quadratic there is
irreducible: a rational-function square root of its
discriminant would be regular at $v=3$, because
the discriminant is regular and nonzero there, and
would specialize to a rational square root of the
nonsquare value (P-F73h).  All other noncollision
conditions follow, for example, by their values
at $v=3$.  Hence the generic central element $G$
has minimal polynomial exactly $\mathcal P(T)$:
each of the five product components is present.
For this specific cubic source example, generic tensor
faithfulness is also verified explicitly in GCT VII,
Section~7.1.3, equations~(25)--(26), printed pp.~21--22:
the actual tensor cube decomposes as
$\bigoplus_{i=0}^5 W_i\otimes T_i$, where the $T_i$
are all the simple source modules over its splitting
field and all $W_i$ are nonzero, of dimensions
$16,4,4,8,2,6$ respectively.  Thus each simple
component occurs in the actual generic action,
which is therefore faithful.  This uses the proved
decomposition for this example, not faithfulness of
an arbitrary abstract nonstandard algebra in a tensor
image.  Its characteristic polynomial
$\chi_G(T)$ is consequently divisible by
$\mathcal P(T)$ over $\mathbb Q(v)$.

Both polynomials have Laurent-polynomial
coefficients, and $\mathcal P$ is monic.
Polynomial long division by this monic polynomial
therefore takes place already in
$\mathbb Q[v,v^{-1}][T]$.  Its remainder is zero
over the fraction field and hence zero in that
Laurent ring.  It follows that
\[
             \mathcal P_q(T)\mid\chi_{G,q}(T)
                               \quad\text{in }\mathbb Q[T].
                                                               \tag{P-F73n}
\]
In the semisimple algebra of (P-F73b), the kernel
of any representation is a direct sum of simple
components.  Omitting one component would omit
its distinct central factor from the minimal and
characteristic polynomials of $G$ on that
representation.  But (P-F73n), together with the
coprimality in step~2, excludes every such omission.
Thus every simple component acts nontrivially,
and the specialized actual tensor representation
is faithful.  This proves the claimed identification
with the full source tensor image at $q=3^e$.

\par\medskip\noindent\textbf{4. Geometric realization and word order.}
The projective cubic
$Y^2Z-X^3+XZ^2=0$ is smooth in characteristic $3$.
On the chart $Z\ne0$, its partial derivative with
respect to $X$ is $Z^2\ne0$.  At infinity its only
point is $[0:1:0]$, where its partial derivative
with respect to $Z$ is nonzero.  A smooth plane
cubic has genus one, and this rational point supplies
the origin.  It defines the elliptic curve (P-F73c)
over $\mathbb F_3$ and hence over $\mathbb F_Q$.

Realize the scalar component on one copy of
$\mathcal E$, the three $M_2(\mathbb Q)$ components
on three pairs $\mathcal E^2$, and $M_2(K)$ on
$\mathcal A^2$ using the faithful inclusion of
P-F70.  All rational scalars act as rational
multiplication endomorphisms.  Product inclusions
and projections make these block diagonal actions
a faithful unital homomorphism into
$\operatorname{End}_{\mathbb F_Q}^0(\mathcal X)$.
This proves (P-F73e); no assertion that the full
endomorphism algebra of $\mathcal X$ is block
diagonal is needed.

Let $\tau$ be identity on the scalar block and
transposition on each matrix block of (P-F73b).
It is an anti-automorphism because all coefficient
fields are commutative.  If $\jmath$ denotes the
block diagonal geometric inclusion, set
\[
                 \rho_\ell(b)=
                   \bigl(\jmath(\tau\Phi(b))\bigr)^*.
                                                               \tag{P-F73o}
\]
Pullback on $H^1$ is additive and reverses
composition; Theorem~15.1(a) of Milne's
\emph{Abelian varieties}, printed p.~27,
identifies it with the dual Tate-module action:
\url{https://www.jmilne.org/math/xnotes/AVs.pdf}.
Thus the two reversals in (P-F73o) make
$\rho_\ell$ a left algebra action, preserving
every ordered source word.  Pullback transposes
the component matrix of a map on a product,
so the resulting action on $H^1$ has exactly
the matrices $M_i(h)$ in (P-F73b), not their
transposes.

More explicitly, put
\[
                  V_{\rm sum}=\mathbb Q
                          \oplus(\mathbb Q^2)^3\oplus K^2.
                                                               \tag{P-F73p}
\]
This is the direct sum of the rational simple
modules in (P-F73b), of rational dimension $11$.
On the elliptic factors $H^1(\mathcal E)$ is
two-dimensional, giving two copies of each
indicated rational module.  On $\mathcal A$,
P-F69 gives the rank-one $L\otimes\mathbb Q_\ell$
module, where $L/K$ is quadratic; choosing a
$K$-basis of $L$ gives two copies of $K^2$ on
$H^1(\mathcal A^2)$.  Consequently, as actual
source modules,
\[
 H^1(\mathcal X,\mathbb Q_\ell)
      \simeq(V_{\rm sum}\otimes\mathbb Q_\ell)^{\oplus2}.
                                                               \tag{P-F73q}
\]
Every simple component occurs, so this action is
faithful, with dimension $22$.  The weight-one
purity is the abelian-variety Weil property,
already used in P-F69.  Scaling either generator
by the actual rational value $F(q)$ commutes
with the construction and preserves its precise
source normalization.
\end{proof}

\paragraph{The remaining geometric comparison.}
This theorem realizes the whole actual specialized
$21$-dimensional source algebra and its ordered
products, extending P-F70's quadratic block.
Its operators are rational quasi-endomorphisms
commuting with Frobenius.  They are not being
identified with Frobenius itself.  The construction
does not preserve the formal $q$-grading before
specialization, assert effective correspondences,
or supply a nonnegative multiplication rule for
the source's candidate canonical basis.  It is
fiberwise in $q=3^e$, with the compatibility
limitations of P-F69 and P-F72 unchanged.
